\documentclass[12pt, twoside]{article}
%\documentclass[12pt, twoside]{article}
\usepackage[letterpaper, margin=1in, headsep=0.2in]{geometry}
\setlength{\headheight}{0.6in}
%\usepackage[english]{babel}
\usepackage[utf8]{inputenc}
\usepackage{microtype}
\usepackage{amsmath}
\usepackage{amssymb}
%\usepackage{amsfonts}
\usepackage[nomessages]{fp} %\FPeval{\var-name}{2*sin(pi/6)}
\usepackage{siunitx} %units in math. eg 20\milli\meter
\usepackage{yhmath} % for arcs, overparenth command
\usepackage{tikz} %graphics
\usetikzlibrary{quotes, angles, arrows, arrows.meta}
\usepackage{pgfplots}
\usepgfplotslibrary{fillbetween}
\pgfplotsset{compat=1.16}
\usepackage{graphicx} %consider setting \graphicspath{{images/}}
\usepackage{parskip} %no paragraph indent
\usepackage{enumitem}
\usepackage{multicol}
\usepackage{venndiagram}

\usepackage{fancyhdr}
\pagestyle{fancy}
\fancyhf{}
\renewcommand{\headrulewidth}{0pt} % disable the underline of the header
\raggedbottom
\hfuzz=2mm %suppresses overfull box warnings

\usepackage{hyperref}

\fancyhead[LO]{La Scuola d'Italia / Dr.\ Huson / IB Math 11 (a) \\* 28 October 2026}
\fancyhead[RO]{\fbox{\begin{minipage}{6cm}\footnotesize Nickname:\\[0.5cm]\end{minipage}}}
\fancyhead[LE]{\fbox{\begin{minipage}{6cm}\footnotesize Nickname:\\[0.5cm]\end{minipage}}}
\fancyfoot[C]{\thepage}

\begin{document}
\subsubsection*{2.6 Quiz: Lines, Exponents and Radicals}

\emph{A calculator is allowed. When a question asks for an exact answer or
a particular form, show the algebra that leads to it.}

\begin{enumerate}[itemsep=0.15cm]

%--- Problem 1: Gradient and equation of AB; perpendicular L through C; intersection F (not the midpoint) [10 marks] ---
%--- Source: Original (AA SL 2.1); parallel of 2-5PQ_Sequences-lines-exponents P3 ---
\item[1.] [Maximum mark: 10]

The points $A(-4,\ 1)$, $B(4,\ 5)$ and $C(1,\ 6)$ are shown on a
coordinate plane.

\begin{enumerate}[itemsep=0.1cm]
  \item Find the gradient of the line $AB$. \hfill [2]
  \item Find the equation of the line $AB$ in the form $y = mx + c$.
  \hfill [2]
\end{enumerate}

The line $L$ passes through $C$ and is perpendicular to the line $AB$.

\begin{enumerate}[itemsep=0.1cm]
  \item[(c)] Find the equation of $L$ in the form $y = mx + c$. \hfill [3]
  \item[(d)] Find the coordinates of the point where $L$ meets the line
  $AB$. \hfill [3]
\end{enumerate}

\begin{tikzpicture}
  \draw (-0.6,0) rectangle (14.6,13.8);
  \foreach \y in {1,2,...,13} \draw[dotted] (0,\y)--(14.1,\y);
\end{tikzpicture}

\newpage

%--- Problem 2: Cinema tickets: set up and solve a 2x2 system; number of children [7 marks] ---
%--- Source: Original (AA SL 2.1, systems of linear equations); parallel of 2-5PQ_Sequences-lines-exponents P4 ---
\item[2.] [Maximum mark: 7]

At a cinema, an adult ticket costs $x$ euros and a child ticket costs $y$
euros. A group of $2$ adults and $3$ children pays $47$ euros. A group of
$1$ adult and $4$ children pays $41$ euros.

\begin{enumerate}[itemsep=0.1cm]
  \item Write down two equations in $x$ and $y$. \hfill [2]
  \item Find the price of an adult ticket and the price of a child ticket.
  \hfill [3]
  \item A family of $2$ adults and some children pays $61$ euros. Find
  the number of children in the family. \hfill [2]
\end{enumerate}

\begin{tikzpicture}
  \draw (-0.6,0) rectangle (14.6,16.4);
  \foreach \y in {1,2,...,15} \draw[dotted] (0,\y)--(14.1,\y);
\end{tikzpicture}

\newpage

%--- Problem 3: Exponent laws; exact fractional negative power; same-base equations [10 marks] ---
%--- Source: Original (AA SL 1.5, 1.7); parallel of 2-5PQ_Sequences-lines-exponents P5 (no equation in u) ---
\item[3.] [Maximum mark: 10]

\begin{enumerate}[itemsep=0.1cm]
  \item Simplify $\left(2a^{3}b\right)^{2} \times 3a^{-2}b^{-4}$. Give your
  answer using positive exponents only. \hfill [3]
  \item Find the exact value of $\left(\dfrac{8}{27}\right)^{-\frac{2}{3}}$.
  \hfill [2]
  \item Solve the equation $9^{x - 1} = 27^{x - 3}$. \hfill [3]
  \item Solve the equation $2^{x + 3} = \dfrac{1}{32}$. \hfill [2]
\end{enumerate}

\begin{tikzpicture}
  \draw (-0.6,0) rectangle (14.6,16.4);
  \foreach \y in {1,2,...,15} \draw[dotted] (0,\y)--(14.1,\y);
\end{tikzpicture}

\newpage

%--- Problem 4: Surds: collect like surds, conjugate product, hence rationalize, linear equation with surd coefficients [8 marks] ---
%--- Source: Original (AA SL 1.7 / prior learning); parallel of 2-5PQ_Sequences-lines-exponents P6 ---
\item[4.] [Maximum mark: 8]

In this question, give every answer in exact form.

\begin{enumerate}[itemsep=0.1cm]
  \item Simplify $\sqrt{50} + \sqrt{8} - \sqrt{32}$. \hfill [2]
  \item Expand and simplify $\left(3 + \sqrt{5}\right)\left(3 - \sqrt{5}\right)$.
  \hfill [2]
  \item Hence, write $\dfrac{8}{3 - \sqrt{5}}$ in the form $a + b\sqrt{5}$,
  where $a, b \in \mathbb{Z}$. \hfill [2]
  \item Solve the equation $x\sqrt{2} + 4 = 3x\sqrt{2}$. Give your answer
  in its simplest form. \hfill [2]
\end{enumerate}

\begin{tikzpicture}
  \draw (-0.6,0) rectangle (14.6,15.9);
  \foreach \y in {1,2,...,15} \draw[dotted] (0,\y)--(14.1,\y);
\end{tikzpicture}

\end{enumerate}
\end{document}
